\documentclass[12pt]{article}

%% Basic document formatting
\usepackage{amsmath, amsthm, amssymb, amsfonts}
\usepackage{mathtools}
\usepackage{xspace}
\usepackage{thmtools}
\usepackage{graphicx}
\usepackage{tikz}
\usepackage{ytableau}
\usepackage{blkarray}
\usepackage{hyperref}
\usepackage{setspace}
\usepackage{array}
\usepackage{fancyhdr}
\usepackage{titling}
\usepackage[left=0.4in,right=0.4in,top=1in,bottom=1in]{geometry}
\usepackage{float}
\usepackage{cancel}
\usepackage{tabularx}
\usepackage[utf8]{inputenc}
\usepackage[english]{babel}
\usepackage{framed}
\usepackage[dvipsnames]{xcolor}
\usepackage{environ}
\usepackage{tcolorbox}
\tcbuselibrary{theorems,skins,breakable}

\usepackage{enumitem}
\setlist[enumerate]{leftmargin=*}
\setlist[enumerate,1]{labelindent=\parindent}
\setlist[enumerate,2]{labelindent=0pt}

% Blackboard Bold
\newcommand{\Z}{\mathbb{Z}}                 % integers
\newcommand{\Zpl}{\mathbb{Z}_{+}}           % positive integers
\newcommand{\Znn}{\mathbb{Z}_{\geq 0}}      % nonnegative integers. 
\newcommand{\Q}{\mathbb{Q}}                 % rationals
\newcommand{\Qpl}{\mathbb{Q}_{+}}           % positive Rationals
\newcommand{\R}{\mathbb{R}}                 % reals
\newcommand{\Rpl}{\mathbb{R}_{+}}           % positive Reals
\newcommand{\C}{\mathbb{C}}                 % complex numbers
\newcommand{\F}{\mathbb{F}}                 % field

% Words
\newcommand{\st}{\text{ such that }}
\newcommand{\wrt}{\text{ with respect to }}
\newcommand{\with}{\text{ with }}
\newcommand{\ie}{\text{, i.e. }}
\newcommand*{\ditto}{\text{---}\texttt{"}\text{---}}

% Operators
\newcommand{\abs}[1]{\left|#1\right|}                   % absolute value
\newcommand{\norm}[1]{\abs{\abs{#1}}}
\newcommand{\floor}[1]{\left\lfloor #1 \right\rfloor}   % floor
\newcommand{\ceil}[1]{\left\lceil #1 \right\rceil}      % ceiling
\newcommand{\powset}[1]{\mathscr{P}(#1)}

% Category Theory 
\renewcommand{\hom}{\text{Hom}}
\newcommand{\homspace}[3]{\hom_{#1}(#2, #3)}
\newcommand{\coproduct}[9]{
  \begin{tikzcd}[ampersand replacement=\&]
      \& \& #1 \arrow[dll, bend right, "#6"] \arrow[dl, "#5"] \\
   #4 \& #3 \arrow[l, "#9"] \\
      \& \& #2 \arrow[ull, bend left, "#8"] \arrow[ul, "#7"]
  \end{tikzcd}
}

\newcommand{\product}[9]{ 
  \begin{tikzcd}[ampersand replacement=\&]
      \& \& #3 \\
      #1 \arrow[r, "#7"] \arrow[urr, bend left, "#5"] \arrow[drr, bend right, "#6"] \& #2 \arrow[ur, "#8"] \arrow[dr, "#9"] \\ 
      \& \& #4
  \end{tikzcd}
}


% Algebra 
\newcommand{\Syl}[2]{\text{Syl}_{#1}{(#2)}}           % set of Sylow-p subgroups 
\newcommand{\<}{\langle}                % \<x\>, subgroup generated by x 
\renewcommand{\>}{\rangle}
\renewcommand{\index}[2]{[#1 : #2]}
\newcommand{\id}{\text{id}}             % identity element
\newcommand{\lhdneq}{%
  \mathrel{\ooalign{$\lneq$\cr\raise.22ex\hbox{$\lhd$}\cr}}}
\newcommand{\order}[1]{\text{o}(#1)}    % order of an element
\newcommand{\spec}{\operatorname{Spec}}
\newcommand{\rad}{\operatorname{rad}}   % radical of an ideal 
\newcommand{\sym}[1]{\mathfrak{S}_{#1}}
\newcommand{\aut}[1]{\text{Aut}(#1)} 
\newcommand{\gal}[2]{\text{Gal}(#1 / #2)} 
\newcommand{\inn}[1]{\text{Inn}(#1)} 
\let\oldcong\cong
\let\oldequiv\equiv
\renewcommand{\cong}{\oldequiv}
\renewcommand{\equiv}{\oldcong}

\newcommand{\triangleleftneq}{\mathrel{\ooalign{%
  $\trianglelefteq$\cr
  \hidewidth\raise0.06ex\hbox{$\not\mkern4mu$}\hidewidth\cr
}}}

\makeatletter % cycle
\newcommand{\cyc}[1]{(\mathbf{\cyc@process#1\relax})}
\def\cyc@process#1#2\relax{%
	#1%
	\ifx\relax#2\relax
	\else
		\,\cyc@process#2\relax
	\fi
}
\makeatother

\newcommand{\GL}[2]{\text{GL}_{#1}(#2)}                             % general linear group
\newcommand{\SL}[2]{\text{SL}_{#1}(#2)}                             % special linear group
\newcommand{\gl}[2]{\mathfrak{gl}_{#1}(#2)}
\renewcommand{\sl}[2]{\mathfrak{sl}_{#1}(#2)}
\newcommand{\so}[2]{\mathfrak{so}_{#1}(#2)}
\newcommand{\su}[1]{\mathfrak{su}(#1)}

% Linear Algebra
\newcommand{\bmat}[1]{\begin{bmatrix}#1\end{bmatrix}}   % bracketed matrix
\newcommand{\rank}{\operatorname{rank}}                 % rank
\newcommand{\nullity}{\operatorname{nullity}}           % nullity
\newcommand{\tr}{\operatorname{tr}}

% Combinatorics
\newcommand{\qnum}[1]{\left[#1\right]_q}                % q-analog
\newcommand{\qfact}[1]{\left[#1\right]!_q}              % q-analog factorial 
\newcommand{\qbinom}[2]{\binom{#1}{#2}_q}               % Gaussian binomial coefficient

% Topology/Analysis
\newcommand{\ball}[2]{\text{B}_{#1}(#2)}  % B_r(x): open r-balls around x
\newcommand{\diam}{\text{diam}}           % diamter of a set in metric space 
\newcommand{\lowint}[2]{\underline{\int_{#1}^{#2}}}
\newcommand{\upint}[2]{\overline{\int}_{#1}^{#2}}

% Misc Notation
\newcommand{\oneton}[1]{\{1, \ldots, #1\}}
\newcommand{\defeq}{\vcentcolon=}   % :=
\newcommand{\eqdef}{=\vcentcolon}   % =: 
\renewcommand{\bf}[1]{\textbf{#1}}
\newcommand{\im}{\operatorname{im}}
\newcommand{\fnrest}[2]{\left.#1\right|_{#2}}
\newcommand{\isom}{\overset{\sim}{\rightarrow}} 


\renewcommand{\thesection}{\S \arabic{section}}
\renewcommand{\thesubsection}{\arabic{subsection}.}
\renewcommand{\thesubsubsection}{(\alph{subsubsection})} 

\newtheorem{claim}{Claim}
\newtheorem*{lemma}{Lemma}

%% Headers & title setup
\newcommand{\course}{\textit{Algebra}, Dummit and Foote} 
\newcommand{\myname}{Jay Ser}
\setlength{\headheight}{14.5pt}
\pagestyle{fancy}
\fancyhf{}
\renewcommand{\headrulewidth}{0.4pt}
\lhead{\course}
\rhead{\myname}
\cfoot{\thepage}
\setlength{\droptitle}{-4em} 
\title{\course\ - Chapter 13: Field Theory}
\author{\myname}
\date{March 25, 2026 $\sim$ \today}

\begin{document}
\maketitle
\thispagestyle{fancy}

%------------------------------------------------------------------------------%
\setcounter{section}{1}
\section{Algebraic Extensions}
\setcounter{subsection}{7} 
\subsection{}
Since the minimal polynomials for $\sqrt{D_1}$ and $\sqrt{D_2}$ are both degree 2, it is clear that regardless of whether $\sqrt{D_1 D_2}$ is a square, $F(\sqrt{D_1}, \sqrt{D_2})$ is a degree 2 or degree 4 extension over $F$. 
Suppose $D_{1} D_{2} = a^2$ for some $a \in F$. 
Since $D_1$ is not a square in $F$, $F(\sqrt{D_{1}})$ is a degree 2 extension over $F$. 
Then $\sqrt{D_{2}}$ satisfies the polynomial $x^2 - D_2$, which can be manipulated as 
\begin{align*}
  x^2 - D_2 & = 0 \\ 
  D_{1} x^2 - D_{1} D_{2} & = 0 \\ 
  x^2 & = \frac{a^2}{D_{1}} \\ 
  x & = \frac{a}{\sqrt{D_1}}
\end{align*}
hence $\sqrt{D_2} \in F(\sqrt{D_1})$, which shows that $F(\sqrt{D_1}, \sqrt{D_2})$ is a degree 2 extension over $F$.  

Conversely, suppose $F(\sqrt{D_1}, \sqrt{D_2})$ is a degree 2 extension over $F$. 
Then necessarily, $\sqrt{D_2} \in F(\sqrt{D_1})$, which means 
\begin{equation} \label{eq:2_8_rtD2_def}
  \sqrt{D_2} = u + v\sqrt{D_1} \hspace{1.5em} (u, v \in F).
\end{equation} 
Square both sides to obtain 
\begin{equation} \label{eq:2_8_sqr} 
  D_2 = u^2 + 2uv\sqrt{D_1} + v^2 D_1.
\end{equation}
Since $\sqrt{D_1} \notin F$, \eqref{eq:2_8_sqr} forces $u = 0$ or $v = 0$. 
By the same reason, \eqref{eq:2_8_rtD2_def} forces $v \neq 0$, so $u = 0$. 
Hence multiplying both sides of \eqref{eq:2_8_sqr} by $D_1$ gives 
$$D_1 D_2 = v^2 D_{1}^{2},$$
i.e., $D_1 D_2$ is a square in $F$, as was to be shown. 
The contrapositive of this result says that if $D_1 D_2$ is not a square, then $F(\sqrt{D_1}, \sqrt{D_2})$ is a degree 4 extension over $F$. 


\subsection{}
Suppose $\sqrt{a + \sqrt{b}} = \sqrt{m} + \sqrt{n}$ for $m, n \in F$. 
Then $a + \sqrt{b} = m + n + 2\sqrt{mn}$ in $F(\sqrt{b})$. 
Since $1, \sqrt{b}$ is a basis for $F(\sqrt{b})$ over $F$ and $\sqrt{b} \notin F$, 
\begin{align*}
  4mn & = b \\
  m + n &  = a 
\end{align*}
where the first equation is from matching $2\sqrt{mn} = \sqrt{b}$. 
Substituting these equations for $a$ and $b$ yields 
$$a^2 - b = m^2 - 2mn + n^2 = (m - n)^2,$$
so $a^2 - b$ is a square in $F$. 

Conversely, suppose $a^2 - b = c^2$ for some $c \in F$. 
As in the forward direction, $m, n \in F$ satisfying the system $4mn = b$ and $m + n = a$ must be found. 
Solving this system for $n$ yields the polynomial 
$$4n^2 - 4an + b = 0$$ 
in $n$. 
Since $\text{char}{F} \neq 2$, the quadratic equation applies to give 
\begin{align*}
  n & = \frac{4a + \sqrt{16a^2 - 16b}}{8} \\ 
    & = \frac{4a + 4c}{8} \\ 
    & = \frac{a + c}{2}.
\end{align*}
It follows $m = \frac{a - c}{2}$. 
By definition, such $m, n \in F$ satisfy $a + \sqrt{b} = m + n + 2\sqrt{mn}$, as desired.  

If $a^2 - b$ is not a square, the result proven above says that $F(\sqrt{a + \sqrt{b}})$ cannot be a degree 2 or a degree 4 extension.
If $a^2 - b$ is a square, then $F(\sqrt{a + \sqrt{b}}) = F(\sqrt{m} + \sqrt{n}) = F(\sqrt{m}, \sqrt{n})$, where the second equality follows by the following observation: 
$F(\sqrt{m} + \sqrt{n}) \subset F(\sqrt{m}, \sqrt{n})$ is obvious. 
The other inclusion follows because $\frac{m - n}{\sqrt{m} + \sqrt{n}} = \sqrt{m} - \sqrt{n}$, where the left side is clearly in $F(\sqrt{m} + \sqrt{n})$. 
By 8), $F(\sqrt{m}, \sqrt{n})$ is biquadratic if and only if $m$, $n$, and $mn$ are not squares in $F$. 
But since $4mn = b$ and $b$ is not a square by assumption, $mn$ is not a square. 
To conclude, $F(\sqrt{a + \sqrt{b}})$ is biquadratic if and only if $a^2 - b$ is a square and one of $m$ or $n$ is not a square. 

\setcounter{subsection}{13}
\subsection{}
$F(\alpha^2) \subset F(\alpha)$ automatically. 
If $F(\alpha) \subsetneq F(\alpha^2)$, then $F(\alpha)$ is a nontrivial extension of $F(\alpha^2)$. 
$\alpha$ satisfies $x^2 - \alpha^2$ over $F(\alpha^2)$, and this is the minimal polynomial for $\alpha$ since the degree of $\alpha$ in $F(\alpha^2)$ must be greater than 1. 
$[F(\alpha) : F] = [F(\alpha) : F(\alpha^2)] [F(\alpha^2) : F]$ says that $2 \mid [F(\alpha) : F]$, which contradicts the assumption that the degree of $\alpha$ over $F$ is odd.  

\subsection{} 
Suppose the claim is not true. 
Let $f(x)$ be minimal with respect to being an irreducible, odd-degree polynomial such that some root $\alpha$ of $f(x)$ gives an extension $F(\alpha)$ of $F$ that is not formally real. 
Let $d \defeq \deg \alpha = \deg f$. 
Since $F(\alpha) \equiv F[x] / \<f(x)\>$, there exist $h_{1}(x), \ldots, h_{t}(x) \in F[x]$ with $\deg{h_i} < d$ such that 
$$-1 = h_{1}(\alpha)^2 + \ldots h_{t}(\alpha)^2$$
in $F(\alpha)$. 
Then there is $g(x) \in F[x]$ such that 
\begin{equation} \label{eq:2_15_fr} 
  -1 + f(x) g(x) = h_{1}(x)^2 + \ldots h_{t}(x)^2
\end{equation}
in $F[x]$. 
$\deg{h_i} < d \implies \deg{h_{i}^2} < 2d \implies \deg{fg} < 2d \implies \deg{g} < d$. 
Furthermore, $\deg g$ is odd: 
without loss of geneerality, suppose $h_{1}, \ldots, h_{k}$ have maximal degree amongst the $h_i$'s. 
Denote their leading terms by $c_i$, respectively. 
Then $\sum_{i = 1}^{k} lt(p_{i}(x)^2) = \sum_{i = 1}^{k} (lt(p_{i}(x)))^2 = \sum_{i = 1}^{k} c_{i}^{2}$, and this sum can't be 0: 
if it was, 
$$c_{2}^{2} + \ldots + c_{k}^{2} = -c_{1}^{2} \implies (c_{2} c_{1}^{-1})^2 + \ldots + (c_{k} c_{1}^{-1})^2 = -1,$$
which contradicts the assumption that $F$ is formally real. 
So the degree of the right side of \eqref{eq:2_15_fr}, $\deg[h_{1}(x)^2 + \ldots + h_{t}(x)^2]$, is even. 
That $\deg{f}$ is odd forces $\deg{g}$ even. 
It follows that $g$ has an irreducible factor of odd degree, say $p(x)$.  
Let $\beta$ be a root of $p$. 
Then in $F(\beta)$, \eqref{eq:2_15_fr} is equivalent to 
$$-1 = h_1(\beta)^2 + \ldots + h_t(\beta)^2,$$ 
hence $F(\beta)$, given by $p(x)$, is not formally real. 
But $p(x)$ is an odd-degree polynomial of a lower degree than $f(x)$, contradicting the minimality of $f$. 

\subsection{} 
Take $\alpha \in R$. 
$\alpha$ is algebraic over $F$, so let $f(x) \in F[x]$ be the minimal polynomial for $\alpha$ over $F$. 
Then the homomorphism 
$$\varphi:F[x] \rightarrow R; \: g(x) \mapsto g(\alpha)$$ 
is a well defined ring homomorphism with $\ker \varphi = \<f(x)\>$, which means $F(\alpha) \subset R$. 
Hence $\alpha^{-1} \in R$, which shows that $R$ is a field.  

\end{document}
